% ============================================================
% rank-fusion (XMTC) — esqueleto ACM acmart/sigconf, pronto para Overleaf.
% Opção [nonacm]: compila um rascunho sem exigir metadados de conferência;
% troque por \documentclass[sigconf]{acmart} + \acmConference{...} na submissão.
% Os \todo{} são pendências reais (ver "## \todo pendentes" na conversa).
% ============================================================
\documentclass[sigconf,nonacm]{acmart}

\usepackage{booktabs}

% \todo inline (robusto no layout duas-colunas). Para notas de margem, troque por:
%   \usepackage[textsize=tiny]{todonotes}
\newcommand{\todo}[1]{\textbf{[TODO: #1]}}

\settopmatter{printacmref=false}        % remove o bloco de referência ACM no rascunho
\setcopyright{none}

\begin{document}

\title{Fusion of Sparse and Dense Rankings for Tail-Label Retrieval in
Extreme Multi-label Text Classification}
% \todo{título de trabalho — refinar antes da submissão}

\author{Daniel Neves Pinheiro}
\affiliation{%
  \institution{Universidade Federal de Minas Gerais (UFMG)}
  \city{Belo Horizonte}
  \country{Brazil}}
\email{} % \todo{e-mail}

\author{Mateus F. Zaparoli Monteiro}
\affiliation{%
  \institution{Universidade Federal de Minas Gerais (UFMG)}
  \city{Belo Horizonte}
  \country{Brazil}}
\email{} % \todo{e-mail}

\begin{abstract}
\todo{Abstract pendente (objetivo -> método -> resultado ancorado -> conclusão;
~150--250 palavras, sem citações). Escrever por último, depois dos Results.}
\end{abstract}

\maketitle

% ============================================================
% INTRODUCTION
% Toda afirmação factual tem âncora % [ficha §N] arquivo:linha (ficha 2026-06-11, commit 1d2f2ef)
% ============================================================
\section{Introduction}

Extreme Multi-label Text Classification (XMTC) assigns to each document a subset of
labels drawn from a space that can reach hundreds of thousands of candidates, whose
frequency distribution is heavily long-tailed: a small \emph{head} of frequent labels
coexists with a vast \emph{tail} of rare ones~\cite{Jain2016PropensityXML, You2019AttentionXML}.
Classifier-based approaches scale to this regime with partitioned label trees or
sparse one-vs-all models~\cite{Prabhu2018Parabel, Babbar2017DiSMEC}, but rare labels
remain the hardest to predict and the easiest to overlook. A complementary view casts
label assignment as information retrieval: the document is treated as a \emph{query}
and the labels are retrieved and ranked by relevance.
% [ficha §1] src/retrieve_sparse.py:5-7 ; src/retrieve_dense.py:4-8 ; CONTEXT.md:8

Under this view, two families of rankers offer complementary signals. A \emph{sparse}
lexical ranker matches surface terms, as in BM25~\cite{robertson2009probabilistic},
while a \emph{dense} ranker compares learned embeddings and captures semantic
similarity~\cite{Karpukhin2020DPR}; the two tend to err in different ways, so combining
their rankings can reinforce correct labels and dampen noise. Rank fusion provides the
machinery for this combination, from score-based operators such as CombSUM and
CombMNZ~\cite{Fox1994CombMNZ} to rank-based methods such as Reciprocal Rank
Fusion~\cite{Cormack2009RRF} and Condorcet fusion~\cite{Montague2002Condorcet}. Because
different rankers emit scores on incomparable scales, a normalization step precedes
fusion. França et al.~\cite{Franca2025RankFusionXMTC} study fusion for XMTC and report
CombMNZ paired with Zero-Mean Unit-Variance normalization (ZMUV) as their strongest
combination.
% [ficha §3] esparso BM25 src/retrieve_sparse.py:64-65,126 ; denso BERT src/retrieve_dense.py:19-23
% [ficha §4] baseline CombMNZ+ZMUV src/fusion.py:5,91-92

Two gaps motivate this work. First, the \emph{interaction} between the fusion
algorithm and the normalization strategy, evaluated specifically for tail labels,
remains under-explored: most reported results aggregate over the whole label space,
where head labels dominate and tail behavior is obscured. Reading the tail requires
either a head/tail split or propensity-scored metrics such as PSP@k and
PSnDCG@k~\cite{Jain2016PropensityXML}. Second, prior fusion tables report only fused
pairs, without isolated single-ranker baselines, which makes the marginal effect of
fusion on the tail difficult to quantify.
% [ficha §4] segmentação head/tail Pareto 0.20 src/metrics.py:53 ; PSP/PSnDCG pendente src/metrics.py:11
% [ficha §6] CSV só tem combos fundidos, sem linha sparse/dense isolada

This paper presents a systematic comparison of ranking-fusion algorithms against
score-normalization strategies for tail-label retrieval in XMTC. We build two base
rankers --- a sparse BM25 k-nearest-neighbor ranker and a fine-tuned BERT bi-encoder
with optional LLM-generated label descriptions --- which we treat as \emph{input},
not as the contribution, generating them once and fusing them offline. Over this pair
we evaluate a grid that currently spans 7 normalizations and 14 fusion algorithms
(98 combinations) and is designed to be extended as further techniques are added. We
adopt 5-fold cross-validation over the pooled dataset, segment Precision@k, nDCG@k and
Recall@k into head and tail (Pareto 80/20), and use CombMNZ+ZMUV as the reference pair.
% [ficha §1] insumo vs contribuição src/fusion.py:16-17 ; CONTEXT.md:20
% [ficha §2] grade aberta 7×14=98 src/fusion.py:40-68,271 ; gridsearch.py:68-69
% [ficha §4] 5-fold CV agrupado src/splits.py:87-109 ; métricas src/metrics.py:51-53

\noindent Our contributions are:
\begin{itemize}
  \item A reproducible, \emph{open} evaluation grid of normalization $\times$ fusion
        combinations for XMTC, currently 7 normalizations $\times$ 14 fusion
        algorithms, with CombMNZ and RRF reimplemented and validated against
        \texttt{ranx}.
        % [ficha §2] src/fusion.py:40-68
  \item A tail-focused evaluation protocol that reports head/tail-segmented P@k,
        nDCG@k and Recall@k under 5-fold cross-validation, averaged as
        mean~$\pm$~standard deviation across folds.
        % [ficha §4] src/metrics.py:53,149-159
  \item An empirical study across XMTC benchmarks of how the choice of normalization
        and fusion algorithm affects tail labels relative to the CombMNZ+ZMUV
        reference. We report results on Eurlex-4K and Wiki10-31K and extend the
        protocol to the larger AmazonCat-13K and Amazon-670K.
        % [ficha §5] eurlex/wiki10 com CSV; amazoncat/670k PENDENTE
        \todo{AmazonCat-13K e Amazon-670K: results/gridsearch.csv pendentes (ficha §5)}
\end{itemize}

\noindent \todo{Frase de achado principal: o ganho da fusão sobre o recuperador
isolado na cauda (documentado como tail nDCG@5 0.45$\to$0.52, +15--17\% vs.\ denso em
CONTEXT.md) NÃO é rastreável a um CSV --- o gridsearch.csv só contém combos fundidos.
Resolver de um destes modos: (a) gerar um CSV com linhas sparse/dense isoladas e citar
o número medido; ou (b) reformular como observação qualitativa sem o percentual. Valor
medido disponível hoje: tail nDCG@5 do par CombMNZ+ZMUV = 0.518 $\pm$ 0.006 no
Eurlex-4K (ficha §6).}

% ============================================================
% Seções seguintes (a escrever via skill escrita-artigo + extrator-de-fatos)
% ============================================================
\section{Related Work}
\todo{Related Work — organizar por tema; âncoras França 2025/SIGIR/SBBD + UDLF
(Valem2023RFE, Pedronette2019LHRR, Valem2018CPRR) e as externas já no refs.bib.}

\section{Method}
\todo{Method — recuperadores base (insumo) + grade fusão × normalização.}

\section{Experimental Setup}
\todo{Setup — datasets, 5-fold CV agrupado, baseline CombMNZ+ZMUV, métricas.}

\section{Results}
\todo{Results — tabelas segmentadas cabeça/cauda (skill /consolidar-resultados).}

\section{Conclusion and Limitations}
\todo{Conclusão + limitações reais (texto stemizado, truncamento 512, 3 de 5 folds,
PSP pendente, denso não salva pesos).}

\bibliographystyle{ACM-Reference-Format}
\bibliography{refs}

\end{document}
